\SGASubsection{7.12}{}
\SGARefTarget{sga3:xxvi:subsec:7:12}

We now study how \(\mathcal R(G/S)\) varies with \(S\).  Let \(S'\)
be an \(S\)-scheme which is likewise nonempty, connected, and
semilocal.  If \(Q\) is a maximal split torus of \(G\), then
\(Q_{S'}\) is a split torus of \(G_{S'}\).  Choose a maximal split
torus \(Q'\) of \(G_{S'}\) containing \(Q_{S'}\), and put
\[
\begin{aligned}
 M&=\operatorname{Hom}_{S\text{-}\mathrm{gr}}
       (Q,\mathbf G_{m,S})
   \simeq
   \operatorname{Hom}_{S'\text{-}\mathrm{gr}}
       (Q_{S'},\mathbf G_{m,S'}),\\
 M'&=\operatorname{Hom}_{S'\text{-}\mathrm{gr}}
       (Q',\mathbf G_{m,S'}).
\end{aligned}
\]
The monomorphism \(Q_{S'}\to Q'\) induces an epimorphism
\[
 u:M'\longrightarrow M.
\]
Write
\[
 L=\Centr_G(Q),\qquad L'=\Centr_{G_{S'}}(Q').
\]
Then \(L'\subset L_{S'}\).

If \(H\) is a subgroup of \(G\) containing \(L\), then \(H_{S'}\)
contains \(L'\), and
\[
\begin{aligned}
 \operatorname{Lie}(H)
  &=\operatorname{Lie}(L)
    \oplus\bigoplus_{\alpha\in R_H}
       \operatorname{Lie}(H)^\alpha,\\
 \operatorname{Lie}(H_{S'})
  &=\operatorname{Lie}(L')
    \oplus\bigoplus_{\alpha'\in R'_{H_{S'}}}
       \operatorname{Lie}(H_{S'})^{\alpha'}.
\end{aligned}
\]
Here \(R_H\), respectively \(R'_{H_{S'}}\), denotes the set of
roots of \(H\), respectively \(H_{S'}\), relative to \(Q\),
respectively \(Q'\).  It follows immediately that
\[
 R_H\subset u\bigl(R'_{H_{S'}}\bigr)
       \subset R_H\cup\{0\}.
\]

Taking \(H=G\), we first obtain
\[
 R\subset u(R')\subset R\cup\{0\}.
\]
Next take for \(H\) a minimal parabolic subgroup \(P\) with Levi
subgroup \(L\).  Then \(R'_{H_{S'}}\) contains a system of positive
roots in \(R'\).  By \SGARefLink{sga3:xxvi:theorem:7-4}{7.4}, there is therefore a minimal parabolic
subgroup \(P'\) of \(G_{S'}\), with Levi subgroup \(L'\), contained
in \(P_{S'}\).  We have constructed the diagram of inclusions
\[
\begin{array}{ccccc}
 Q_{S'} & \subset & L_{S'} & \subset & P_{S'}\\[1.2ex]
 \cap   &         & \cup   &         & \cup\\[1.2ex]
 Q'     & \subset & L'     & \subset & P'.
\end{array}
\]

If \(R^+\), respectively \(\Delta\), is the system of positive,
respectively simple, roots in \(R\) defined by \(P\), and if
\(R'^+,\Delta'\) are defined similarly, then
\[
 R^+\subset u(R'^+)\subset R^+\cup\{0\},
 \qquad
 \Delta\subset u(\Delta')\subset\Delta\cup\{0\}.
\]

Now let
\[
 w\in W\simeq\Norm_G(Q)(S)/\Centr_G(Q)(S),
\]
and choose a representative \(n\in\Norm_G(Q)(S)\).  Since
\(\Int(n)Q=Q\), one has \(\Int(n)L=L\), and hence
\(\Int(n)L_{S'}=L_{S'}\).  The tori \(Q'\) and \(\Int(n)Q'\) are
therefore two maximal split tori of \(L_{S'}\), so they are
conjugate by an element \(x\in L(S')\).  We may choose \(x\) so that
\[
 \Int(nx)Q'=Q'.
\]
Thus \(n'=nx\) belongs to \(\Norm_{G_{S'}}(Q')(S')\).  Let \(w'\)
be the image of \(n'\) in
\[
 W'\simeq
 \Norm_{G_{S'}}(Q')(S')/\Centr_{G_{S'}}(Q')(S').
\]
The action of \(w'\) on \(M'\) is compatible with
\(u:M'\to M\), and the induced action on \(M\) is the action of
\(w\).

Using the definition of relative root data together with the
conjugacy theorems gives the following result.

\begin{sgastatement}{\SGAInlineRefTarget{sga3:xxvi:theorem:7-13}Theorem 7.13}
Let \(S,S'\) be nonempty connected semilocal schemes, let
\(S'\to S\) be a morphism of schemes, let \(G\) be an
\(S\)-reductive group, and write
\[
\begin{aligned}
 \mathcal R(G/S)
  &=(\underline M,\underline M^*,\underline R,
     \underline R^*,\underline R^+),\\
 \mathcal R(G_{S'}/S')
  &=(\underline M',\underline M'^*,\underline R',
     \underline R'^*,\underline R'^+)
\end{aligned}
\]
for the based relative root data.  There is a canonical
homomorphism
\[
 u:\underline M'\longrightarrow\underline M
\]
with the following properties:
\begin{enumerate}[label=\textup{(\roman*)}]
\item The homomorphism \(u\) is surjective.
\item For every \(w\in W\), there is an element \(w'\in W'\)
compatible with \(u\) and inducing \(w\) on \(\underline M\).
\item For every subset \(X\subset\underline M\), put
\[
 X^\wedge=X\cap(\underline M\setminus\{0\}).
\]
Then
\[
 u(\underline R'^+)^\wedge=\underline R^+,
 \qquad
 u(\underline\Delta')^\wedge=\underline\Delta.
\]
\item For every parabolic subgroup \(H\), consider
\[
 t_r(H)\subset\underline\Delta,
 \qquad
 t_r(H_{S'})\subset\underline\Delta'.
\]
Then
\[
\begin{aligned}
 t_r(H_{S'})
 &=u^{-1}\bigl(t_r(H)\cup\{0\}\bigr)
       \cap\underline\Delta'\\
 &=\left\{\alpha'\in\underline\Delta'
       \ \middle|\ 
       u(\alpha')\in t_r(H)\text{ or }u(\alpha')=0
   \right\}.
\end{aligned}
\]
\end{enumerate}
\end{sgastatement}
\SGARefTarget{sga3:xxvi:thm:7:13}

\begin{sgastatement}{Remark 7.14}
If \(G\) is splittable over \(S\), its maximal split tori are
maximal tori, and the relative notions introduced here coincide
with the previously defined absolute notions.  The preceding
theorem therefore describes the relative root datum
\(\mathcal R(G/S)\) and the relative type \(t_r\) in terms of the
absolute root datum and absolute type of \(G_{S'}\), where \(S'\) is
chosen so that \(G_{S'}\) is splittable
\textup{(cf.\ \SGARefLink{sga3:xxiv:theorem:4-4-1}{Exposé XXIV, 4.4.1})}.  See \textup{[BT65]}, \SGARefLink{sga3:xxvi:cor:6:11}{6.12}
and following, for this description.
\end{sgastatement}
\SGARefTarget{sga3:xxvi:rem:7:14}

\SGASubsection{7.15}{}
\SGARefTarget{sga3:xxvi:subsec:7:15}

Let \(S\) be a Henselian local scheme, \(s_0\) its closed point, and
\(S_0\) the spectrum of the residue field at \(s_0\), identified
with a closed subscheme of \(S\).  For every object \(X\) over
\(S\), write \(X_0\) for its base change to \(S_0\).  Let \(G\) be
an \(S\)-reductive group.

For every parabolic subgroup \(\overline P\) of \(G\), the group
\(\overline P_0\) is a parabolic subgroup of \(G_0\).  Conversely,
for every parabolic subgroup \(P\) of \(G_0\), there is a parabolic
subgroup \(\overline P\) of \(G\) such that
\[
 \overline P_0=P.
\]
This follows from Hensel's lemma and from the fact that \(\Par(G)\)
is a smooth \(S\)-scheme.  In particular, by \SGARefLink{sga3:xxvi:corollary:5-7}{5.7},
\(\overline P\) is minimal if and only if \(\overline P_0\) is
minimal.

Once such a subgroup \(\overline P\) has been chosen, the same
argument shows that every maximal split subtorus of
\(\overline P_0\) is of the form \(T_0\), where \(T\) is a maximal
split subtorus of \(\overline P\).  It follows readily that the
relative root data of \(G\) over \(S\) and of \(G_0\) over \(S_0\)
are canonically isomorphic.  Thus the theory of parabolic subgroups
of \(G\) reduces to the corresponding theory for \(G_0\).

Conversely, every \(S_0\)-reductive group is of the form \(G_0\), by
\SGARefLink{sga3:xxiv:1.21}{Exposé XXIV, Proposition 1.21}.  This makes it possible to reduce the
study of the parabolic subgroups of an \(S_0\)-reductive group to
the corresponding study over \(S\).

\chapter*{Bibliography}
\addcontentsline{toc}{chapter}{Bibliography}
\SGARefTarget{sga3:xxvi:bibliography}

\begin{description}[leftmargin=4.5em,style=nextline]
\item[\textup{[BT65]}]
A.~Borel and J.~Tits, \emph{Groupes réductifs},
Publ.\ Math.\ I.H.É.S.\ \textbf{27} \textup{(1965)}, 55--150.
\footnote{\emph{Editor's note.} This reference appeared in the
original; the references which follow were added in the
Polo--Gille edition.}

\item[\textup{[BT71]}]
A.~Borel and J.~Tits,\\
\emph{Éléments unipotents et sous-groupes
paraboliques de groupes réductifs. I},
Invent.\ Math.\ \textbf{12} \textup{(1971)}, 95--104.

\item[\textup{[Ch05]}]
C.~Chevalley, \emph{Classification des groupes algébriques
semi-simples}, with the collaboration of P.~Cartier,
A.~Grothendieck, and M.~Lazard, Collected Works, vol.~3,
Springer, 2005.

\item[\textup{[DG70]}]
M.~Demazure and P.~Gabriel, \emph{Groupes algébriques},
Masson \& North-Holland, 1970.

\item[\textup{[Gi71]}]
J.~Giraud, \emph{Cohomologie non abélienne},
Springer-Verlag, 1971.
\end{description}
